\scanpage{458}\mspage{370}
% archive sheet number: 457
\uline{4.14} : \uline{Critères d'accessibilité pour grosses catégories}.

\uline{Théorème} 4.14.1. Soient $C$ une catégorie équivalente à une petite catégorie, $\Phi,\inserted{\Psi}$ \deleted{deux} petits ens de \deleted{\unsure{foncteurs}} catégories, \deleted{\unsure{$\varphi:I\longrightarrow C$} Pour tte} équivalentes à des petites catégories. Pour tte catégorie $I$, on désigne par $I_\bullet$ (resp $ {}_\bullet I$) le cône droit (resp gauche) sur $I$, déduit de $I$ en \unsure{ajoutant} un objet final \unsure{$e_I$} (resp initial \inserted{$\deleted{\unsure{e}}e=e_I$}). Soit $\Lambda_1$ \inserted{(resp. $\Lambda_2$)} un ens. de couples $(\deleted{\unsure{\varphi}}I,\varphi)$ (resp. $(I,\psi)$), où $I\in\Phi$ (resp. $I\in\Psi$), et où
\begingroup\csname tagsleft@true\endcsname\[
\varphi:I_{\bullet}\longrightarrow C\qquad(\text{resp. }\psi:{}_\bullet I\longrightarrow C)
\tag*{(4.14.1.1)}
\]\endgroup
est un foncteur. Pour tte catégorie $\unsure{\mathcal N|N|\mathscr N}$ \unsure{stable}$^*$ par $\unsure{\varinjlim|\underline{\lim}}$ de type $\Phi$ et par $\varprojlim$ de type $\Psi$, on désigne par
% slanted
\margin{Cette propriété de stabilité n'est pas \unsure{néc.} pour la déf. \deleted{\unsure{naïve}} de $\underline{\mathrm{Hom}}_{\Lambda_1}^{\Lambda_2}(C,N)$, et il est \unsure{à peu près|toujours} sûr que \unsure{le th.} est \unsure{valable|applicable} sans faire cette hyp sur $N$. Cf 4.14.16, où \unsure{l'on|on ?} voit \unsure{que|ce} \unsure{pouvant|pour} \unsure{s'étendre|l'établir} \unsure{dans|en} \unsure{toute sa généralité|toute généralité} pour $\Psi$.}
\begingroup\csname tagsleft@true\endcsname\[
\underline{\mathrm{Hom}}_{\Lambda_1}^{\Lambda_2}(C,\mathcal N)
\subset\underline{\mathrm{Hom}}(C,\mathcal N)
\tag*{(4.14.1.2)}
\]\endgroup
la sous-catégorie \unsure{str} pleine de $\underline{\mathrm{Hom}}(C,\mathcal N)$ formée des foncteurs $f:C\longrightarrow\mathcal N$ \deleted{tel} ayant la propriété suivante :

a) $\forall\,(I,\varphi)\in\Lambda_1$, le foncteur composé
\diagram{\begin{tikzcd}[ampersand replacement=\&]
I_\bullet\arrow[r,"\varphi"] \& C\arrow[r,"f"] \& \mathcal N
\end{tikzcd}}
est $I$-\unsure{exact|ex.ct|coex.|coexact} \inserted{\unsure{à g.|à G}}, i.e. fait de $(f\varphi)(e_I)$ une $\underset{I}{\varinjlim}$ de $f\varphi|I$.
